% 阋神星（综述）
% license CCBYSA3
% type Wiki

本文根据 CC-BY-SA 协议转载翻译自维基百科\href{https://en.wikipedia.org/wiki/Eris_(dwarf_planet)}{相关文章}。





















\subsection{另见}
\begin{itemize}
\item 天体命名
\item 清除邻近的小天体
\item 国际天文学联合会（IAU）
\item 假设的海王星外行星
\end{itemize}
\subsection{参考文献}
\begin{enumerate}
\item Staff. Discovery Circumstances: Numbered Minor Planets. IAU: Minor Planet Center. 2007-05-01 [2007-05-05]. （原始内容存档于2016-05-04）.
\item Brown, Mike. The discovery of 2003 UB313 Eris, the largest known dwarf planet. 2006 [2007-05-03]. （原始内容存档于2014-09-10）.
\item JPL Small-Body Database Browser: 136199 Eris (2003 UB313). 2009-11-20 [2009-01-21]. （原始内容存档于2019-01-09）.
\item List Of Centaurs and Scattered-Disk Objects. Minor Planet Center. [2008-09-10]. （原始内容存档于2011-07-25）.
\item Buie, Marc. Orbit Fit and Astrometric record for 136199. Deep Ecliptic Survey. 2007-11-06 [2007-12-08]. （原始内容存档于2014-09-26）.
\item Asteroid Observing Services （页面存档备份，存于互联网档案馆）
\item Sicardy, B.; Ortiz, J. L.; Assafin, M.; Jehin, E.; Maury, A.; Lellouch, E.; Gil-Hutton, R.; Braga-Ribas, F.; Colas, F.; Widemann. Size, density, albedo and atmosphere limit of dwarf planet Eris from a stellar occultation (PDF). European Planetary Science Congress Abstracts. 2011, 6: 137 [September 14, 2011]. Bibcode:2011epsc.conf..137S. （原始内容存档 (PDF)于October 18, 2011）.
\item Beatty, Kelly. Former 'tenth planet' may be smaller than Pluto. NewScientist.com. Sky and Telescope. November 2010 [October 17, 2011]. （原始内容存档于February 23, 2012）.
\item 掩星觀測顯示鬩神星和冥王星大小幾乎相等. PanSci 泛科学. 2011-11-03 [2017-03-09]. （原始内容存档于2014-03-25） （中文（台湾））.
\item Michael E. Brown and Emily L. Schaller. The Mass of Dwarf Planet Eris (abstract page). Science. 2007, 316 (5831): 1585 [2011-02-01]. PMID 17569855. doi:10.1126/science.1139415. （原始内容存档于2010-06-26）.
\item Kelly Beatty. Eris Gets Dwarfed (Is Pluto Bigger?). Sky & Telescope (News Blog). 2010-11-07 [2010-11-07]. （原始内容存档于2011-08-30）.
\item Snodgrass, Carry, Dumas, Hainaut. Characterisation of candidate members of (136108) Haumea's family (abstract). The Astrophysical Journal. 2009-12-16. arXiv:0912.3171.
\item AstDys (136199) Eris Ephemerides. Department of Mathematics, University of Pisa, Italy. [2009-03-16]. （原始内容存档于2011-06-04）.
\item F. Bertoldi, W. Altenhoff, A. Weiss, K.M. Menten, C. Thum. The trans-neptunian object UB313 is larger than Pluto. Nature. 2006-02, 439 (7076): 563–564 [2018-04-02]. ISSN 1476-4687. doi:10.1038/nature04494. （原始内容存档于2019-07-13） （英语）.
\item “阋”，《说文解字》卷三：“恒讼也。《诗》云：‘兄弟阋于墙。’从斗从儿。儿，善讼者也。许激切。”繁：“鬩”，简：“阋”，拼音：xì，注音：ㄒㄧˋ，粤拼：jik1
\item JPL/NASA. What is a Dwarf Planet?. Jet Propulsion Laboratory. 2015-04-22 [2022-01-19]. （原始内容存档于2021-12-08）.
\item How Big Is Pluto? New Horizons Settles Decades-Long Debate. www.nasa.gov. 2015 [2015-07-14]. （原始内容存档于2017-07-01）.
\item Mike Brown. The discovery of 2003 UB313 Eris, the 10th planet largest known dwarf planet. Caltech. 2006 [2010-01-05]. （原始内容存档于2012-05-20）.
\item The IAU draft definition of "planet" and "plutons" (新闻稿). IAU. 2006-08-16 [2006-08-16]. （原始内容存档于2006-08-20）.
\item Mike Brown. The shadowy hand of Eris. Mike Brown's Planets. 2010 [2010-11-07]. （原始内容存档于2010-11-11）.
\item Mike Brown. How big is Pluto, anyway?. Mike Brown's Planets. 2010-11-22 [2010-11-23]. （原始内容存档于2015-10-31）. (Franck Marchis on 2010-11-08)
\item Brown, Mike. How big is Pluto, anyway?. Mike Brown's Planets. 2010-11-22 [2010-11-23]. （原始内容存档于2015-10-31）. (Franck Marchis on 2010-11-08)
\item Thomas H. Maugh II and John Johnson Jr. His Stellar Discovery Is Eclipsed. Los Angeles Times. 2005-10-16 [2008-07-14]. （原始内容存档于2017-02-21）.
\item KLC. 矮行星Eris中譯名正式命名為「鬩神星」. tamweb.tam.gov.tw. 台北天文馆. 2007-07-10 [2017-03-09]. （原始内容存档于2013-07-29） （中文（台湾））.
\item 方舟子. 方舟子：欺世盗名的八卦宇宙论. 人民网. 2005-08-01 [2013-03-30]. （原始内容存档于2013-05-09） （中文（中国大陆））.
\item Gemini Observatory Shows That"10th Planet" Has a Pluto-Like Surface. Gemini Observatory. 2005 [2007-05-03]. （原始内容存档于2007-03-11）.
\item M. E. Brown, C. A.Trujillo, D. L. Rabinowitz. Discovery of a Planetary-sized Object in the Scattered Kuiper Belt. The Astrophysical Journal (abstract page). 2005, 635 (1): L97–L100. Bibcode:2005ApJ...635L..97B. arXiv:astro-ph/0508633 可免费查阅. doi:10.1086/499336.
\item J. Licandro, W. M. Grundy, N. Pinilla-Alonso, P. Leisy. Visible spectroscopy of 2003 UB313: evidence for N2 ice on the surface of the largest TNO (PDF). Astronomy and Astrophysics. 2006, 458 (1): L5–L8 [2011-05-03]. Bibcode:2006A&A...458L...5L. arXiv:astro-ph/0608044 可免费查阅. doi:10.1051/0004-6361:20066028. （原始内容存档 (PDF)于2008-11-21）.
\item M. E. Brown, M. A. van Dam, A. H. Bouchez, D. LeMignant, C. A. Trujillo, R. Campbell, J. Chin, Conrad A., S. Hartman, E. Johansson, R. Lafon, D. L. Rabinowitz, P. Stomski, D. Summers, P. L. Wizinowich. Satellites of the largest Kuiper belt objects. The Astrophysical Journal (abstract page). 2006, 639 (1): L43–L46. Bibcode:2006ApJ...639L..43B. arXiv:astro-ph/0510029 可免费查阅. doi:10.1086/501524.
\end{enumerate}


\subsection{外部链接}

Michael Brown 关于 Eris 和 Dysnomia 的网页  
MPC 关于 (136199) Eris 的数据库条目  
NASA JPL 的 3D 轨道可视化  
2003 UB313 轨道模拟  
Keck 天文台关于 Dysnomia 发现的页面  
AstDyS-2（小行星动力学网站）中的 Eris（矮行星）  
历表 · 观测预测 · 轨道信息 · 本征要素 · 观测信息  
JPL 小天体数据库中的 Eris（可在 Wikidata 编辑）  
近距离接近 · 发现 · 历表 · 轨道查看器 · 轨道参数 · 物理参数













